\documentclass[10pt]{article}
\usepackage[margin=1in]{geometry}
\usepackage{booktabs}
\usepackage{longtable}
\usepackage{hyperref}
\usepackage{listings}
\usepackage{xcolor}
\usepackage{amsmath}
\usepackage{amssymb}
\usepackage{textcomp}

\lstset{basicstyle=\ttfamily\footnotesize,breaklines=true,frame=single,backgroundcolor=\color{gray!8}}
\title{Independent Replication Report: BVBRC-114 \\ \large Khajanchi et al.\ (2017) --- IncFIB Plasmid-Encoded Virulence Factors in Food-Source \emph{Salmonella enterica}}
\author{Ollie (OpenClaw Wave-Keeper Subagent) \\ \small Compute host: uicgpu; Judge model: argo:claude-opus-4.7 (free)}
\date{2026-07-05}

\begin{document}
\maketitle

\begin{abstract}
This report documents an independent computational replication of Khajanchi et al.\ (BMC Genomics 2017, DOI: 10.1186/s12864-017-3954-5; PMID: 28768482; PMC5541697), covering all four \emph{in silico}--testable core claims: (i) serotype composition (6 \emph{S.} Typhimurium + 1 \emph{S.} Heidelberg among the 7 focal isolates); (ii) IncFIB plasmid presence in the 6 Typhimurium strains and absence in the SE819 recipient; (iii) 99--100\% amino-acid conservation of plasmid-encoded SitABCD and aerobactin (IucABC) across the 6 Typhimurium isolates versus the CP001122.1 (pCVM29188\_146) IncFIB(K) reference; and (iv) the 5+1 monophyletic-subclade topology (5 Typhimurium tight cluster; SE397 outlier). All four \emph{in silico} claims are reproduced on the actual public WGS data under NCBI BioProject \textbf{PRJNA312617}. Two additional wet-lab claims (Caco-2 persistence of the SE819::IncFIB transconjugant; iron-responsive qRT-PCR of \emph{sit}/\emph{iuc}) are explicitly not testable in an \emph{in silico} replication and are marked as such. \textbf{Verdict: PARTIAL} in the strong sense (bioinformatic core reproduced; wet-lab claims out of computational reach).
\end{abstract}

\section{Paper summary}
The paper sequences seven \emph{S. enterica} food/animal-source isolates (six \emph{S.} Typhimurium---SE163A, SE397, SE452, SE478, SE696A, SE710A---and one \emph{S.} Heidelberg SE819), compares them against $\sim$45 publicly-available \emph{S.} Typhimurium/Heidelberg genomes, and characterises IncFIB-plasmid-encoded iron-acquisition operons (SitABCD, aerobactin \emph{iucABCD-iutA}). Its main conclusions: (i) 6+1 serovar breakdown reproduces from SeqSero; (ii) 5 of 6 Typhimurium strains form a monophyletic subclade, 1 branches with bovine-source Typhimurium; (iii) Sit and aerobactin operons on the IncFIB plasmid are highly conserved across bacterial species; (iv) transferring the IncFIB plasmid into IncFIB-deficient SE819 yields a transconjugant that persists in Caco-2 cells at a higher rate than the recipient; (v) \emph{sit}/\emph{iuc} genes are iron-repressible (upregulated under iron limitation).

\section{Claims table}
\begin{longtable}{@{}p{0.5cm}p{6cm}p{2cm}p{2.5cm}p{1cm}p{3cm}@{}}
\toprule
\# & Claim & Type & Testable? & Tested? & Outcome \\ \midrule \endhead
C1 & 7 WGS assemblies public under PRJNA312617 & Data availability & YES & \checkmark & REPLICATED (all 7 downloaded) \\
C2 & 6 Typhimurium + 1 Heidelberg (SE819) & In-silico serotype & YES & \checkmark & \textbf{REPLICATED exactly} (SeqSero2 + MLST orthogonal) \\
C3 & 5-strain monophyletic Typhimurium subclade; 1 outlier & Phylogeny & YES & \checkmark & REPLICATED (direction) --- SE397 is the outlier \\
C4 & 6 Typhimurium carry IncFIB; SE819 does not & Plasmidome & YES & \checkmark & \textbf{REPLICATED exactly} (PlasmidFinder DB blastn) \\
C5 & Sit + aerobactin operons conserved & Sequence conservation & YES & \checkmark & REPLICATED --- 99.65--100\% AA identity \\
C6 & SE819::IncFIB persists in Caco-2 & Wet-lab virulence & NO & --- & NOT TESTABLE in silico \\
C7 & \emph{sit}/\emph{iuc} iron-responsive & Wet-lab qRT-PCR & NO & --- & NOT TESTABLE in silico \\ \bottomrule
\end{longtable}

\section{Method (numbered)}
\subsection{Assembly retrieval}
Seven focal assemblies were retrieved: SE163A (GCA\_001647955.1), SE696A (GCF\_001700345.1), SE710A (GCF\_001700365.1), SE819 (GCF\_001647935.1) via \texttt{datasets download genome accession}; SE397 (GCF\_001729025.1), SE452 (GCF\_001729035.1), SE478 (GCF\_001729045.1) via direct HTTPS FTP fetch (\texttt{ftp.ncbi.nlm.nih.gov/genomes/all/GCF/001/729/...}) as a fallback around an intermittent \texttt{api.ncbi.nlm.nih.gov} DNS failure on uicgpu. Reference plasmid \textbf{CP001122.1} (pCVM29188\_146; 146,811 bp; IncFIB(K)+IncFII; explicitly cited in paper Table~1) plus LT2 (NC\_003197.1), CVM29188 chromosome (CP001121.1), and three bovine Typhimurium (BovineChina SA972816 CP007484.1, USMARC-1808 CP014969.1, USMARC-1880 CP014981.1) via NCBI E-utilities.

\subsection{In-silico serotyping (C2)}
\begin{lstlisting}[language=bash]
SeqSero2_package.py -m k -t 4 -i <assembly.fna> -d <outdir> -p 8
\end{lstlisting}
Results in \texttt{evidence/seqsero/seqsero\_summary.tsv}.

\subsection{Orthogonal MLST}
\begin{lstlisting}[language=bash]
mlst --scheme salmonella <assembly.fna>   # tseemann/mlst v2.35.0
\end{lstlisting}
Results in \texttt{evidence/mlst/mlst\_auto.tsv}. 6/6 Typhimurium $\rightarrow$ identical ST19; SE819 $\rightarrow$ ST15 (canonical Heidelberg ST).

\subsection{PlasmidFinder-style Inc-rep detection (C4)}
Cloned CGE PlasmidFinder DB (488 rep sequences). For each assembly:
\begin{lstlisting}[language=bash]
makeblastdb -in <asm.fna> -dbtype nucl -out <db>
blastn -query plasmid_all.fsa -db <db> -evalue 1e-20 -perc_identity 80 \
       -outfmt "6 qseqid sseqid pident length qlen sstart send evalue"
# Post-filter: coverage = length/qlen >= 60%
\end{lstlisting}

\subsection{Iron-acquisition operon tblastn (C4/C5)}
CDS translations for \emph{sitA/B/C/D + iucA/B/C + iutA + iroB} extracted from CP001122.1 GenBank record via Biopython + regex on \texttt{/gene=} and \texttt{/product=}. For each of the 7 focal genomes:
\begin{lstlisting}[language=bash]
tblastn -query iron_proteins.faa -db <db> -evalue 1e-20 -max_target_seqs 5
# Presence: %ID >= 90 AND coverage >= 90% of query
\end{lstlisting}

\subsection{Phylogeny (C3)}
\texttt{mash sketch -k 21 -s 1000} on 12 genomes (7 focal + 5 refs); all-vs-all \texttt{mash dist}; Biopython \texttt{DistanceTreeConstructor.nj} for the NJ tree.

\section{Results vs paper (section-by-section)}

\subsection{4.1 Serotype (C2) --- \textbf{REPLICATED EXACTLY}}
6/7 assemblies serotyped as \emph{S.} Typhimurium (antigenic profile 1,2), SE819 as \emph{S.} Heidelberg. MLST orthogonal: 6/6 Typhimurium = identical ST19, SE819 = ST15. What worked: SeqSero2 k-mer mode was fast ($\sim$15--30~s per assembly), well-calibrated, and matched MLST perfectly. What didn't: no failures.

\subsection{4.2 IncFIB plasmid presence (C4) --- \textbf{REPLICATED EXACTLY}}
6/6 Typhimurium strains carry \texttt{IncFIB(AP001918)\_1} at 98.09\% identity, full 682/682 coverage. SE819 has NO IncFIB rep. Additional Inc replicons detected across strains: IncFIA, IncFII variants, IncA, IncX4, IncI1, IncC, ColRNAI, ColpVC. What worked: PlasmidFinder DB is authoritative, blastn is deterministic. What didn't: the \emph{extra} replicons detected in each strain were not addressed in the paper (see Open Question Q5).

\subsection{4.3 Iron acquisition operon presence and conservation (C4/C5) --- \textbf{REPLICATED}}
\begin{center}
\begin{tabular}{lccccccc} \toprule
Gene & SE163A & SE397 & SE452 & SE478 & SE696A & SE710A & SE819 \\ \midrule
sitA & YES (99.7) & YES (100) & YES (99.7) & YES (99.7) & YES (99.7) & YES (99.7) & NO (76.1 chr) \\
sitB & YES (100) & YES (100) & YES (100) & YES (100) & YES (100) & YES (100) & NO (77.7 chr) \\
sitC & YES (99.7) & YES (99.7) & YES (99.7) & YES (99.7) & YES (99.7) & YES (99.7) & NO (85.8 chr) \\
sitD & YES (100) & YES (100) & YES (100) & YES (100) & YES (100) & YES (100) & NO (68.5 chr) \\
iucA & YES (100) & YES (100) & YES (100) & YES (100) & YES (100) & YES (100) & --- \\
iucB & YES (99.7) & YES (99.7) & YES (99.7) & YES (99.7) & YES (99.7) & YES (99.7) & --- \\
iucC & YES (100) & YES (100) & YES (100) & YES (100) & YES (100) & YES (100) & --- \\ \bottomrule
\end{tabular}
\end{center}
What worked: 99.65--100\% AA identity across 6 independently isolated strains (3+ US states, 1992--2002) is striking direct confirmation of the paper's conservation claim. What didn't: (a) iucD, iroC, iroN, shfB, pefA, ompX were not annotated on CP001122.1 with matching \texttt{/gene=} fields, so we did not query them; (b) the 6 Typhimurium plasmids are fragmented across $\sim$3 assembly contigs each ($\sim$15--20 kb sit+iuc contig separate from the IncFIB rep contig), so a naive ``same contig as rep'' check is misleading (see Open Question Q1).

\subsection{4.4 Phylogeny (C3) --- \textbf{REPLICATED (direction)}}
Mean mash distance from each Typhimurium focal to the other 5:
\begin{center}
\begin{tabular}{lc} \toprule
Strain & Mean $d \times 10^{-3}$ \\ \midrule
SE163A & 1.71 \\
SE452 & 1.73 \\
SE478 & 1.81 \\
SE696A & 1.48 \\
SE710A & 1.51 \\
\textbf{SE397} & \textbf{4.29} $\rightarrow$ outlier \\ \bottomrule
\end{tabular}
\end{center}
What worked: SE397 is unambiguously the outlier at 2.5$\times$ higher within-group distance --- matching paper Fig~1b's ``5+1'' topology. What didn't: mash resolution is limited; SE397 in our NJ tree does not cleanly resolve as sister to bovine references specifically (paper's core-SNP tree does). Closing that gap would require Parsnp/RAxML on core-SNP alignment ($\sim$1~h additional uicgpu time; deferred).

\subsection{4.5 Wet-lab claims C6 (Caco-2 persistence) and C7 (iron-responsive qRT-PCR)}
Not testable in silico. Reported as ``not tested'' rather than failed or contradicted. Would require BSL-2 wet-lab access.

\section{Verdict + justification}
\textbf{PARTIAL REPLICATION (strong sense).} Rationale:
\begin{itemize}
\item 4/4 testable in-silico core claims independently reproduced at high confidence.
\item 2 wet-lab claims explicitly out of scope for a computational replication.
\item Not full \textsc{REPLICATED} because C6+C7 cannot be confirmed here.
\item Not \textsc{SPOT-CHECK} because we ran actual analyses on actual data, not merely verified availability.
\end{itemize}

\section{Open Questions}
See \texttt{open\_questions.json} for machine-readable form.
\begin{description}
\item[Q1.] Do our 6 assemblies systematically break the $\sim$140 kb IncFIB plasmid at a repeat-mediated boundary? \emph{Next step:} long-read resequencing of one strain to close the plasmid and identify the break point.
\item[Q2.] Is SE397 genuinely sister to bovine-source Typhimurium at core-SNP resolution, and does it carry bovine-associated accessory genes? \emph{Next step:} Parsnp + RAxML core-SNP tree + panaroo pangenome.
\item[Q3.] Is the 99.65--100\% conservation of plasmid \emph{sit}+\emph{iuc} driven by ongoing purifying selection or a recent single-source HGT? \emph{Next step:} PAML codeml $dN/dS$ on plasmid vs chromosomal \emph{sit} alignments.
\item[Q4.] Which plasmid factor (\emph{sit}/\emph{iuc}/\emph{pef}/\emph{spv}) carries the causal Caco-2-persistence signal? \emph{Next step:} SalCom + Kroger 2013 intracellular RNA-seq cross-reference; then wet-lab $\Delta$-loci constructs.
\item[Q5.] Are the additional Inc/Col replicons in the 6 strains co-integrates with IncFIB or independent smaller plasmids that could confound transconjugant attribution? \emph{Next step:} long-read closed assembly of the SE819::IncFIB transconjugant.
\end{description}

\section*{Data and code availability}
All input assemblies are public under NCBI BioProject PRJNA312617. This replication's scripts, evidence tables, and log are in \\ \texttt{\textasciitilde/Dropbox/REPLICATE-PROJECT/BVBRC-114-Salmonella-incompatibility-Khajanchi2017/}.
\end{document}
